% GENERATED FILE. Edit canonical lessons, then run npm run generate:books.
% canonical-source-hash: fnv1a64:e445339c494aec5a
% canonical-lessons: ZH-C329-tonghua, ZH-C329-gongzhu, ZH-C329-xiannu, ZH-C329-xiaoshi, ZH-C329-jieju, ZH-R329-review-324-329

\chapter{Once upon a Time}
\label{ch:zh-once-upon-a-time}
\hlchaptermodality{329}

\begin{chapteropening}
\textbf{By the end of this chapter:} \emph{I can talk about fairy tales, a princess and a fairy, say someone vanished, and talk about how a story ends.}

Complete the last lesson of chapter 329: I can talk about fairy tales, a princess and a fairy, say someone vanished, and talk about how a story ends.
\end{chapteropening}

\section[tónghuà]{tónghuà — fairy tale}

\label{lesson:ZH-C329-tonghua}



\begin{quote}
Before the new one: say the Mandarin for a dragon, then the Mandarin for to carry messages between people.
\end{quote}

\subsection*{You'll want to know: tónghuà}
\emph{tónghuà} — \textquotedblleft{}fairy tale\textquotedblright{}.

\begin{quote}
\emph{xiǎoshíhou māma měitiān wǎnshang gěi wǒ dú tónghuà} — When I was little, my mum read me a fairy tale every evening.
\end{quote}

Its characters come later; for now it is a word for the ear.

\begin{cousinweb}
\emph{tónghuà} is built from \emph{tóng} 'a child' and \zh{话} 'talk', both inherited from Middle Chinese, \emph{duwng} and \emph{hwaejH} (Guangyun, in Baxter's transcription), which says of the first that it means a child not yet married.
\end{cousinweb}

\begin{grammarlens}[title={What you've built}]
One more word for telling what happened.
\end{grammarlens}

\subsection*{Your turn}
\begin{itemize}
\raggedright
  \item \emph{Say it:} \emph{tónghuà}
  \item \emph{Say it:} \emph{tónghuà}, once more
  \item \emph{Say it:} \emph{xiǎoshíhou māma měitiān wǎnshang gěi wǒ dú tónghuà} — When I was little, my mum read me a fairy tale every evening.
  \item \emph{Recall:} say the Mandarin for a dragon, then the Mandarin for to carry messages between people, then say \emph{tónghuà} again
\end{itemize}

\begin{tcolorbox}[breakable,colback=teal!4,colframe=teal!35!black,title={Before you move on}]
What does tónghuà mean? (\textquotedblleft{}Fairy tale\textquotedblright{}.) Say it once more.
\end{tcolorbox}
\section[gōngzhǔ]{gōngzhǔ — princess}

\label{lesson:ZH-C329-gongzhu}



\begin{quote}
Before the new one: say the Mandarin for to carry messages between people, then the Mandarin for a fairy tale for children.
\end{quote}

\subsection*{You'll want to know: gōngzhǔ}
\emph{gōngzhǔ} — \textquotedblleft{}princess\textquotedblright{}.

\begin{quote}
\emph{cóngqián yǒu yí ge hěn piàoliang de gōngzhǔ} — Once upon a time there was a very beautiful princess.
\end{quote}

Its characters come later; for now it is a word for the ear.

\begin{cousinweb}
\emph{gōngzhǔ} is built from \zh{公} and \emph{zhǔ}, both inherited from Middle Chinese, \emph{kuwng} and \emph{tsyuX} (Guangyun, in Baxter's transcription), which glosses the first 'a father' and 'an official' and the second 'to be in charge' and 'a lord'. How the pair came to mean a princess is not settled.
\end{cousinweb}

\begin{grammarlens}[title={What you've built}]
One more word for telling what happened.
\end{grammarlens}

\subsection*{Your turn}
\begin{itemize}
\raggedright
  \item \emph{Say it:} \emph{gōngzhǔ}
  \item \emph{Say it:} \emph{gōngzhǔ}, once more
  \item \emph{Say it:} \emph{cóngqián yǒu yí ge hěn piàoliang de gōngzhǔ} — Once upon a time there was a very beautiful princess.
  \item \emph{Recall:} say the Mandarin for to carry messages between people, then the Mandarin for a fairy tale for children, then say \emph{gōngzhǔ} again
\end{itemize}

\begin{tcolorbox}[breakable,colback=teal!4,colframe=teal!35!black,title={Before you move on}]
What does gōngzhǔ mean? (\textquotedblleft{}Princess\textquotedblright{}.) Say it once more.
\end{tcolorbox}
\section[xiānnǚ]{xiānnǚ — fairy; heavenly maiden}

\label{lesson:ZH-C329-xiannu}



\begin{quote}
Before the new one: say the Mandarin for a fairy tale for children, then the Mandarin for a princess in a story.
\end{quote}

\subsection*{You'll want to know: xiānnǚ}
\emph{xiānnǚ} — \textquotedblleft{}fairy; heavenly maiden\textquotedblright{}.

\begin{quote}
\emph{chuánshuō lǐ yǒu qī ge xiānnǚ} — In the legend there are seven fairies.
\end{quote}

Its characters come later; for now it is a word for the ear.

\begin{cousinweb}
\emph{xiānnǚ} is built from \emph{xiān} 'an immortal' and \zh{女} 'a woman', both inherited from Middle Chinese, \emph{sjen} and \emph{nrjoX} (Guangyun, in Baxter's transcription), which says of the first that one who grows old and does not die is called so.
\end{cousinweb}

\begin{grammarlens}[title={What you've built}]
One more word for telling what happened.
\end{grammarlens}

\subsection*{Your turn}
\begin{itemize}
\raggedright
  \item \emph{Say it:} \emph{xiānnǚ}
  \item \emph{Say it:} \emph{xiānnǚ}, once more
  \item \emph{Say it:} \emph{chuánshuō lǐ yǒu qī ge xiānnǚ} — In the legend there are seven fairies.
  \item \emph{Recall:} say the Mandarin for a fairy tale for children, then the Mandarin for a princess in a story, then say \emph{xiānnǚ} again
\end{itemize}

\begin{tcolorbox}[breakable,colback=teal!4,colframe=teal!35!black,title={Before you move on}]
What does xiānnǚ mean? (\textquotedblleft{}Fairy; heavenly maiden\textquotedblright{}.) Say it once more.
\end{tcolorbox}
\section[xiāoshī]{xiāoshī — to disappear; to vanish}

\label{lesson:ZH-C329-xiaoshi}



\begin{quote}
Before the new one: say the Mandarin for a princess in a story, then the Mandarin for a heavenly maiden.
\end{quote}

\subsection*{You'll want to know: xiāoshī}
\emph{xiāoshī} — \textquotedblleft{}to disappear; to vanish\textquotedblright{}.

\begin{quote}
\emph{dì'èr tiān zǎoshang nàge xiānnǚ tūrán xiāoshī le} — The next morning the fairy suddenly vanished.
\end{quote}

Its characters come later; for now it is a word for the ear.

\begin{cousinweb}
\emph{xiāoshī} is built from \emph{xiāo} and \emph{shī}, both inherited from Middle Chinese, \emph{sjew} and \emph{syit} (Guangyun, in Baxter's transcription), which glosses the first 'to die out' and 'to be used up', and the second 'to err' and 'to let go'.
\end{cousinweb}

\begin{grammarlens}[title={What you've built}]
One more word for telling what happened.
\end{grammarlens}

\subsection*{Your turn}
\begin{itemize}
\raggedright
  \item \emph{Say it:} \emph{xiāoshī}
  \item \emph{Say it:} \emph{xiāoshī}, once more
  \item \emph{Say it:} \emph{dì'èr tiān zǎoshang nàge xiānnǚ tūrán xiāoshī le} — The next morning the fairy suddenly vanished.
  \item \emph{Recall:} say the Mandarin for a princess in a story, then the Mandarin for a heavenly maiden, then say \emph{xiāoshī} again
\end{itemize}

\begin{tcolorbox}[breakable,colback=teal!4,colframe=teal!35!black,title={Before you move on}]
What does xiāoshī mean? (\textquotedblleft{}To disappear; to vanish\textquotedblright{}.) Say it once more.
\end{tcolorbox}
\section[jiéjú]{jiéjú — ending}

\label{lesson:ZH-C329-jieju}



\begin{quote}
Before the new one: say the Mandarin for a heavenly maiden, then the Mandarin for to vanish without a trace.
\end{quote}

\subsection*{You'll want to know: jiéjú}
\emph{jiéjú} — \textquotedblleft{}ending\textquotedblright{}.

\begin{quote}
\emph{kàn le jiéjú yǐhòu wǒ kū le} — I cried after reading the ending.
\end{quote}

Its characters come later; for now it is a word for the ear.

\begin{cousinweb}
\emph{jiéjú} is built from \emph{jié} 'to tie' and \emph{jú}, both inherited from Middle Chinese, \emph{ket} and \emph{gjowk} (Guangyun, in Baxter's transcription), which glosses the first 'to tie' and the second 'an office' and 'a portion'. How the pair came to mean an ending is not settled.
\end{cousinweb}

\begin{grammarlens}[title={What you've built}]
One more word for telling what happened.
\end{grammarlens}

\subsection*{Your turn}
\begin{itemize}
\raggedright
  \item \emph{Say it:} \emph{jiéjú}
  \item \emph{Say it:} \emph{jiéjú}, once more
  \item \emph{Say it:} \emph{kàn le jiéjú yǐhòu wǒ kū le} — I cried after reading the ending.
  \item \emph{Recall:} say the Mandarin for a heavenly maiden, then the Mandarin for to vanish without a trace, then say \emph{jiéjú} again
\end{itemize}

\begin{tcolorbox}[breakable,colback=teal!4,colframe=teal!35!black,title={Before you move on}]
What does jiéjú mean? (\textquotedblleft{}Ending\textquotedblright{}.) Say it once more.
\end{tcolorbox}
\section[(dialogue)]{Looking back: from Remembering to Once upon a Time}

\label{lesson:ZH-R329-review-324-329}



\begin{quote}
Say the Mandarin for to call to mind and to honour the memory of.
\end{quote}

\subsection*{Your turn}
\begin{itemize}
\raggedright
  \item \emph{Say it:} to call to mind, to honour the memory of, to leave behind for others, an age in history, an impression left on you
  \item \emph{Say it:} to give an account of, to pass away, one's descendants, for generations on end, to miss someone dearly
  \item \emph{Say it:} a nation with its own government, the whole wide world, the outside world, to found something new, an ancient civilization
  \item \emph{Say it:} a famous person, a hero, an emperor, to go into battle, a war between countries
  \item \emph{Say it:} a myth about the gods, an old folk tale, to be handed down over time, a dragon, to carry messages between people
  \item \emph{Say it:} a fairy tale for children, a princess in a story, a heavenly maiden, to vanish without a trace, the ending of a story
  \item \emph{Read it:} \zh{世} \zh{界} \zh{传} — each one aloud, then a word that holds it
\end{itemize}

\begin{tcolorbox}[breakable,colback=teal!4,colframe=teal!35!black,title={Before you move on}]
To call to mind? (\textbf{\zh{想起}}, \emph{xiǎngqǐ}.) To honour the memory of? (\emph{jìniàn}.) To leave behind for others? (\emph{liúxià}.) An age in history? (\textbf{\zh{时代}}, \emph{shídài}.) An impression left on you? (\emph{yìnxiàng}.) To give an account of? (\emph{jiǎngshù}.) To pass away? (\textbf{\zh{去世}}, \emph{qùshì}.) One's descendants? (\textbf{\zh{后代}}, \emph{hòudài}.) For generations on end? (\textbf{\zh{世代}}, \emph{shìdài}.) To miss someone dearly? (\emph{xiǎngniàn}.) A nation with its own government? (\textbf{\zh{国家}}, \emph{guójiā}.) The whole wide world? (\textbf{\zh{世界}}, \emph{shìjiè}.) The outside world? (\textbf{\zh{外界}}, \emph{wàijiè}.) To found something new? (\emph{jiànlì}.) An ancient civilization? (\textbf{\zh{文明}}, \emph{wénmíng}.) A famous person? (\textbf{\zh{名人}}, \emph{míngrén}.) A hero? (\emph{yīngxióng}.) An emperor? (\emph{huángdì}.) To go into battle? (\emph{dǎzhàng}.) A war between countries? (\emph{zhànzhēng}.) A myth about the gods? (\emph{shénhuà}.) An old folk tale? (\textbf{\zh{传说}}, \emph{chuánshuō}.) To be handed down over time? (\emph{liúchuán}.) A dragon? (\emph{lóng}.) To carry messages between people? (\textbf{\zh{传话}}, \emph{chuánhuà}.) A fairy tale for children? (\emph{tónghuà}.) A princess in a story? (\emph{gōngzhǔ}.) A heavenly maiden? (\emph{xiānnǚ}.) To vanish without a trace? (\emph{xiāoshī}.) The ending of a story? (\emph{jiéjú}.)
\end{tcolorbox}
